\documentclass[12pt]{article}

% --------------------
% Paquetes
% --------------------
\usepackage[spanish]{babel}
\usepackage[utf8]{inputenc}
\usepackage[T1]{fontenc}
\usepackage{setspace}
\usepackage{geometry}
\usepackage{times}
\usepackage{graphicx}
\usepackage{amsmath}
\usepackage{soul}
\usepackage{tabularx}
\usepackage{booktabs}

% --------------------
% Márgenes
% --------------------
\geometry{
  letterpaper,
  left=2.5cm,
  right=2.5cm,
  top=2.5cm,
  bottom=2.5cm
}

\geometry{margin=2.5cm} % Ajusta los márgenes para tener más espacio
\usepackage{listings}


% Configuración de listings para que el código no se salga de la página
\lstset{
    basicstyle=\ttfamily\small,
    breaklines=true,         % Activa el salto de línea automático
    breakatwhitespace=false, % Permite saltos aunque no haya espacios
    columns=fullflexible,
    frame=single,            % Añade un recuadro opcional
    keepspaces=true,
    literate={á}{{\'a}}1 {é}{{\'e}}1 {í}{{\i}}1 {ó}{{\'o}}1 {ú}{{\'u}}1 {ñ}{{\~n}}1 % Soporte para acentos
}

% --------------------
% Documento
% --------------------
\begin{document}
\thispagestyle{empty}

% --------------------
% Portada
% --------------------
\begin{center}

{\LARGE \textbf{Instituto Politécnico Nacional}}\\[0.5cm]
{\large Escuela Superior de Cómputo}\\[1.5cm]



{\Huge \textbf{Procesamiento de Lenguaje Natural}}\\[0.5cm]
{\large Tercer Parcial}\\[2cm]

{\LARGE \textbf{Tareas}}\\[0.5cm]

\end{center}

\vfill

\begin{minipage}{0.55\textwidth}
\textbf{Alumno:}\\
Sánchez García Miguel Alexander\\[1cm]

\textbf{Grupo:} 6AV1\\
\textbf{Carrera:} Licenciatura en Ciencia de Datos
\end{minipage}
\hfill
\begin{minipage}{0.35\textwidth}
\textbf{Profesor:}\\
Ortiz Castillo Marco Antonio\\[1cm]

\textbf{Fecha:}\\
9 de enero, 2026
\end{minipage}

\newpage

% --------------------
% Contenido
% --------------------
\setstretch{1.3}

\section*{\hl{N - gramas}}

Son utilizados para obtener secuencias de ``n'' palabras. La idea es obtener una mayor cantidad de contexto en función de la ``n''.

\vspace{1em}

\textbf{Por ejemplo:} \underline{No} \underline{sé} diferencia entre algoritmos PLN

\vspace{1em}

\begin{minipage}[t]{0.45\textwidth}
    \textbf{Una secuencia para $n=1$}
    \begin{itemize}
        \item No\_sé
        \item sé\_diferencia
        \item diferencia\_entre
        \item \dots
        \item algoritmos\_PLN
    \end{itemize}
\end{minipage}
\hfill
\begin{minipage}[t]{0.45\textwidth}
    \textbf{Para $n=2$}
    \begin{itemize}
        \item No sé \_ diferencia
        \item sé diferencia \_ entre
        \item diferencia entre \_ algoritmos
    \end{itemize}
\end{minipage}

\vspace{2em}

El siguiente paso consiste en vectorizar el conjunto de palabras, con lo cual se genera embeddings.

\vspace{1.5em}

$\rightarrow$ \textbf{PCA}

\begin{itemize}
    \item No [ .5 \quad .3 \quad .9 ]
    \item Sé [ .1 \quad .1 \quad .7 ] , algoritmo de afinación
\end{itemize}


\section*{\hl{Topic Modeling}}

Conjunto de técnicas para descubrir temas ``ocultos'' en una colección de textos, encontrando patrones de co-ocurrencia de palabras. Existe LDA, LSA y LSI.

\subsection*{Latent Dirichlet Allocation (LDA)}
LDA es de los más usados al ser un modelo probabilístico generativo.

\begin{itemize}
    \item Cada doc. se modela como una mezcla de tópicos ($\theta_d$).
    \item Cada tópico se modela como una distribución de palabras ($\rho_k$).
\end{itemize}

\textbf{Hiperparámetros:}
\begin{itemize}
    \item \textbf{$\alpha$:} Influye en la cantidad de tópicos que tiene un doc (pocos dominantes o muchos mezclados).
    \item \textbf{$\beta$:} Influye en si un tópico usa pocas palabras características o una distribución más plana.
\end{itemize}

Para inferir tópicos el modelo suele usar métodos como Bayes variacional, MCMC o Gibbs sampling.

\textbf{Salida:} Proporciones de tópicos por documento y palabras más probables por tópico.

\vspace{1em}

\section*{Notebook: LDA (Latent Dirichlet Allocation)}

\textbf{Descripción:} Este notebook implementa el análisis de tópicos usando Latent Dirichlet Allocation (LDA) con documentos en español. Se realiza preprocesamiento de texto incluyendo tokenización, eliminación de stopwords y lematización antes de aplicar el modelo LDA.

\subsection*{Celdas de Importación y Datos}

\begin{lstlisting}[language=Python]
from gensim.test.utils import common_corpus
import gensim
from gensim.models.ldamulticore import LdaMulticore
from nltk.stem import WordNetLemmatizer
from nltk import word_tokenize
from nltk.corpus import stopwords
import string
import matplotlib.pyplot as plt
from gensim.utils import simple_preprocess
from gensim.models import CoherenceModel
import nltk
nltk.download('stopwords')
nltk.download('wordnet')
nltk.download('omw-1.4')
from nltk.corpus import stopwords
from nltk.stem.wordnet import WordNetLemmatizer
from gensim import corpora
from gensim.models import LdaModel
from gensim.models import lsimodel

# Documentos en ingles
doc_1_en = "Sugar is bad to consume. My sister likes to have sugar, but not my father."
doc_2_en = "The theory of relativity was proposed by Albert Einstein in 1905."
doc_3_en = "Many people believe that sugar is not good for your health."
doc_4_en = "Einstein developed the theory of general relativity."
doc_5_en = "The health benefits of regular exercise are well documented."
doc_6_en = "Climate change is affecting ecosystems around the world."
doc_7_en = "Machine learning algorithms can predict patterns in data."
doc_8_en = "Healthy eating habits include consuming less sugar and more vegetables."
doc_9_en = "Artificial intelligence is revolutionizing many industries."
doc_10_en = "Physical activity reduces the risk of chronic diseases."

# Documentos en espanol
doc_1_es = "El azucar es malo para consumir. A mi hermana le gusta el azucar, pero no a mi padre."
doc_2_es = "La teoria de la relatividad fue propuesta por Albert Einstein in 1905."
doc_3_es = "Mucha gente cree que el azucar no es bueno para la salud."
doc_4_es = "Einstein desarrollo la teoria de la relatividad general."
doc_5_es = "Los beneficios para la salud del ejercicio regular estan bien documentados."
doc_6_es = "El cambio climatico esta afectando los ecosistemas alrededor del mundo."
doc_7_es = "Los algoritmos de aprendizaje automatico pueden predecir patrones en los datos."
doc_8_es = "Los habitos alimenticios saludables incluyen consumir menos azucar y mas vegetales."
doc_9_es = "La inteligencia artificial esta revolucionando muchas industrias."
doc_10_es = "La actividad fisica reduce el riesgo de enfermedades cronicas."
\end{lstlisting}

\subsection*{Procesamiento y Modelado}

\begin{lstlisting}[language=Python]
# Usar documentos en espanol
corpus = [doc_1_es, doc_2_es, doc_3_es, doc_4_es, doc_5_es, doc_6_es, doc_7_es, doc_8_es, doc_9_es, doc_10_es]
stop = set(stopwords.words('spanish'))
exclude = set(string.punctuation)
lemma = WordNetLemmatizer()

def clean(doc):
    stop_free = " ".join([i for i in doc.lower().split() if i not in stop])
    punc_free = ''.join(ch for ch in stop_free if ch not in exclude)
    normalized = " ".join(lemma.lemmatize(word) for word in punc_free.split())
    return normalized

clean_corpus = [clean(doc).split() for doc in corpus]
dictionary = corpora.Dictionary(clean_corpus)
doc_term_matrix = [dictionary.doc2bow(doc) for doc in clean_corpus]

# Entrenar el modelo LDA
lda_model = LdaModel(doc_term_matrix, num_topics=3, id2word=dictionary, passes=15)

# Ver los topicos
print("TOPICOS ENCONTRADOS")
for idx, topic in lda_model.print_topics(-1):
    print(f'\nTopico {idx}:')
    print(topic)

print("DOCUMENTOS PROCESADOS")
print(f"Total de documentos: {len(clean_corpus)}")
print(f"Vocabulario Unico: {len(dictionary)}")
\end{lstlisting}


\end{document}